\section{Random sparse tables and their interfaces}
\label{sec:atom}

The exact hierarchy uses matrix entries, correction products, and column
detector expectations.  We represent each of these by sample-dependent
tables: bins represent entries of the transitions and residuals,
correction tables represent the pilot and actual corrections, and score
tables represent the detector trials.  Each normalized table has an
earlier preparatory table that supplies its denominator.

We first prove the exact conditional-mean identities.  We then extend the
same row formulas to computed raw arrays and prove the support and access
bounds that hold even when those arrays are inaccurate.  Finally we bound
the effect of errors in the incoming cutoffs and column multipliers.  The
later implementation supplies enumeration and pairing of the resulting
ports.

Exact conditional-mean identities refer only to the true tables defined
here.  Applying the same formulas to approximate raw signals gives
approximate tables; their normalization need not be unbiased.

Weighted pseudorandomness provides a related way to represent expectations
with few random bits; see Braverman, Cohen and Garg~\cite{BravermanCohenGarg2020}
and the error-reduction work of Cohen, Doron, Renard, Sberlo and
Ta-Shma~\cite{CohenDoronRenardSberloTaShma2021}. The tables below have
exact conditional-mean identities. Their estimates use a common
environment, so products and subsequent normalization are analyzed
through joint error bounds rather than independent sampling.

\subsection{Table domains and the normalization principle}

Throughout this section $s=8$.  Fix a stage, write $f=b_l$, and let $G$ be
one of its nonnegative raw matrix inputs.  The hierarchy gives absolute
constants $M_G,K_G$ such that
\begin{equation}
 \sum_y G(x,y)\le M_G,\qquad G(x,y)\le K_G f.
 \label{eq:atom-input-bounds}
\end{equation}
These bounds include pilot and actual inputs and the residual $D$; see
Lemmas~\ref{lem:hier-pilot-bounds} and~\ref{lem:hier-residual}.
For a correction $F=fQ(G/f)g_0$, the additional input is
\begin{equation}
 g_0(y)>0\quad\Longrightarrow\quad
 \#\{x:G(x,y)\ge\delta_1 f\}\le R_*/f,
 \label{eq:atom-column-input}
\end{equation}
from Lemma~\ref{lem:hier-column-test}.

A row task has start rate $u\le f$, end scale $a$, and, when needed, a
usage rate $t\le a$, all on the rate grid.  Its domain is
$X_u(x)=Y_a(y)=1$, with distinct start and terminal sampling matrices.  We use
$a\le K_f f$ for an absolute $K_f$, and use $a=f$ for corrections and
detector scores.  Here and in the estimation and compression sections,
$t$ denotes a sampling rate; the numerical digit position $\nu$ is introduced
separately in Section~\ref{sec:arith}.  The true first raw signal is
\begin{equation}
 \mathsf{d}_y=G(x,y)/a.
 \label{eq:atom-first-signal}
\end{equation}
An optional second raw signal is
$\zeta\mathbf1_{\{Y_t(y)=1\}}q_y$, where $q_y=q(x,y)\in[0,1]$ is a
deterministic conditional mean and $\zeta>0$ is a fixed sufficiently small
constant.  That mean will be specified separately for each normalized task.

All parameters of a table are common to the whole table: its stage and
iteration, task type, sampling matrix roles, sampling matrix padding suffixes, rates
$u,a,t$, and subsequently its budget and numerical convention.  They do
not depend on the row or column at which the table is queried.

As in Section~\ref{sec:hierarchy}, $\operatorname{largest}_k$ denotes the
$k$th largest value, with unlimited zero padding.  For two arrays on the same padded index set,
\begin{equation}
 |\operatorname{largest}_k(v)-\operatorname{largest}_k(w)|
 \le\sup_z|v_z-w_z|.
 \label{eq:atom-rank-lipschitz}
\end{equation}
Indeed, if every coordinate changes by at most $\epsilon$, at least $k$
coordinates at or above the old $k$th value remain at or above that value
minus $\epsilon$; apply this in both directions.

The normalization can be seen before choosing the row functions.  Fix a
pair $(x,y)$ on its conditioning slice, a deterministic nonnegative value
$k(\mathsf{d}_y)$, and a random row cutoff $R\in[0,1]$.  Choose a preparatory factor
$h$ that is one wherever $k$ is positive.  If
\[
 q(x,y)=\mathbb E[h(\mathsf{d}_y)R\mid X_u(x)=Y_t(y)=1]\ge\tfrac12
\]
on that support, then $k(\mathsf{d}_y)R/q(x,y)$ has conditional mean $k(\mathsf{d}_y)$
there.  A clipped denominator keeps the formula bounded at every other
position and for computed inputs.  The constructions below arrange a
uniform positive lower bound on this mean while the row cutoff limits the
number of nonzero retained entries.  Corrections use an additional incoming
cutoff, with the same normalization principle.

For real $\alpha<\beta$, let $\chi_{\alpha,\beta}$ be a fixed increasing
switch, zero on $(-\infty,\alpha]$ and one on $[\beta,\infty)$.
Choose it by rescaling the normalized integral of
$z^{s-1}(1-z)^{s-1}$ on $[0,1]$.
Both $\chi_{\alpha,\beta}^{1/s}$ and
$(1-\chi_{\alpha,\beta})^{1/s}$ are Lipschitz: at either endpoint the
vanishing function has a zero of order $s$, and on the interior ordinary
differentiation applies.  Products of bounded functions with Lipschitz
$s$th roots again have Lipschitz $s$th roots.

Every scalar saturation below uses fixed min/max clipping,
\[
 \operatorname{clip}_{[\alpha,\beta]}(v)
   =\min\{\beta,\max\{\alpha,v\}\},
\]
which is $1$-Lipschitz.  We specify its bounds when defining the function.
Functions initially defined for nonnegative first entries are extended
by zero below zero; each already vanishes below a positive support
threshold.  These conventions apply only to evaluation of the row
functions.  Order statistics and access tests always use the full prescribed raw
arrays before scalar clipping.

Here is one explicit set of margins used below.  Let $\vartheta>0$ be a
fixed lower support threshold of a main row function, and put
$\rho=\operatorname{largest}_K(v)$ for its first raw row $v$.  Define
\begin{align}
 h_{\vartheta}(v)
   &=\chi_{\vartheta/2,\,3\vartheta/4}(v),\notag\\
 R_{\vartheta}(\rho)
   &=1-\chi_{\vartheta/1000,\,\vartheta/500}(\rho),\notag\\
 I_{\vartheta}(v,\rho)
   &=\chi_{\vartheta/8,\,\vartheta/4}(v)
       \bigl(1-\chi_{\vartheta/100,\,\vartheta/50}(\rho)\bigr).
 \label{eq:atom-threshold-template}
\end{align}
We take $\vartheta=\theta_1$ for corrections, $\vartheta=1/2$ for 
bins, and $\vartheta=\delta_0$ for the detector score.  For that detector
we make the allowed hierarchy choices
\[
 \eta_-=\delta_0/1000,\qquad \eta_+=\delta_0/500,\qquad
 \rho(t)=1-\chi_{\delta_0/1000,\,\delta_0/500}(t)
          =R_{\delta_0}(t).
\]
Here $\rho$ denotes the cutoff in the hierarchy, and the same switch
$\chi$ is used in both occurrences.  These choices satisfy
$0<\eta_-<\eta_+<\delta_0/3$.  Thus the detector row cutoff here is exactly
$R_i$ in Equation~\eqref{eq:hier-pilot-row-gate}.
The hierarchy's choices may and do satisfy
$\delta_1<\theta_1/100$ and $0<\delta_0<\delta_1$.
Thus the correction auxiliary score $I_{\theta_1}$ can be positive only
when $G(x,y)>\delta_1 f$.

\subsection{Bins and preparatory normalization}

We first split a raw entry into contributions at the rate-grid scales.
At each scale a row cutoff bounds the number of retained entries, and an
exact preparatory conditional mean corrects the loss caused by that cutoff.

Let $\chi=\chi_{1/2,1}$.  For $w\ge0$ define
\begin{align}
 \omega_0(w)&=\chi(w),\notag\\
 \omega_k(w)&=\chi(w/b_k)-\chi(w/b_{k-1})\qquad(k\ge1).
 \label{eq:atom-bin-weights}
\end{align}
The top bin has no preceding switch.  The weights are nonnegative and
sum to one for $w>0$, since their partial sums telescope to
$\chi(w/b_k)\to1$.  Hence
\begin{equation}
 w=\sum_{k\ge0}w\omega_k(w).
 \label{eq:atom-bin-partition}
\end{equation}
At scale $a=b_k$, the bin in units of the sampling scale is
\begin{equation}
 k_a(d)=d\omega_k(ad).
 \label{eq:atom-bin-function}
\end{equation}
For $k\ge1$ its support is contained in $1/2<d<2^H$.
Its two switch transition intervals are disjoint, and
\[
 \chi(d)-\chi(d/2^H)=\chi(d)(1-\chi(d/2^H)),
\]
which proves directly that its $s$th root is Lipschitz.
For $k=0$, fix $B_{\rm top}>\max\{K_G,1\}$ and evaluate the defining
expression at $\operatorname{clip}_{[0,B_{\rm top}]}(d)$.
This makes the top-bin function bounded with Lipschitz root.  It
preserves every true value, since $a=b_0=1$ and
$G(x,y)\le K_G f\le K_G$.  For $k\ge1$, retain the zero extension
outside the support $(1/2,2^H)$; no upper clipping changes a zero bin
value.  Bins with $a>2K_G f$ vanish on the true input, so a fixed
$K_f\ge2K_G$ suffices for the permitted scale range.

For the routed residual bin use instead
\begin{equation}
 k_a^{(j)}(d)=k_a(d)s_j(ad/f).
 \label{eq:atom-routed-bin}
\end{equation}
The routing fractions have bounded Lipschitz roots.  Since $a/f\le K_f$
and the bin is zero above a fixed argument range (or is top-bin clipped),
these functions have uniformly bounded Lipschitz roots as well.  The
same zero-extension and clipping conventions apply to routed bins.
Summing Equation~\eqref{eq:atom-routed-bin} over scales gives the routed
residual contribution exactly on every true input.

Choose $R_{1/2}$ and $h=h_{1/2}$ from
Equation~\eqref{eq:atom-threshold-template}.  For the true row set
\[
 \rho_x=\operatorname{largest}_K
   \{G(x,z)/a:Y_a(z)=1\},\qquad R=R_{1/2}(\rho_x).
\]
Write
\begin{equation}
 q_y=\mathbb E\bigl[h(\mathsf{d}_y)R\mid X_u(x)=Y_t(y)=1\bigr],
 \qquad
 c(q)=\min\{1,\max\{1/2,q\}\}.
 \label{eq:atom-dense-denominator}
\end{equation}
The number $q_y$ is the deterministic conditional mean at the indicated
pair.  The second raw signal carries the $Y_t$ membership mask.  In
addition, $q_y=0$ if $\mathsf{d}_y\le1/4$.
The true preparatory table on the broad
terminal domain $Y_a$ is the unnormalized value $h(\mathsf{d}_y)R$.  It uses only
the first raw signal.  Its conditional mean on the narrower slice is
exactly the $q_y$ in Equation~\eqref{eq:atom-dense-denominator}.
The usage rate $t$ is part of the fixed task tuple: changing the slice
on which this mean is taken would in general change the denominator.
The broad preparatory domain and the narrower normalization slice must
therefore be retained as separate parts of the interface.

\begin{lemma}[Bin table identity]
\label{lem:atom-dense}
Fix an admissible tuple $(u,a,t)$ with $a=b_k$ as above.
For sufficiently large fixed $K$, the true table
\begin{equation}
 \mathcal B_a(x,y)=
  \mathbf1_{\{Y_t(y)=1\}}\frac{k_a(\mathsf{d}_y)R}{c(q_y)}
 \label{eq:atom-dense-table}
\end{equation}
has uniformly bounded values and row support, and is zero outside
$X_u(x)=Y_t(y)=1$. On this slice,
\[
 a\,\mathbb E[\mathcal B_a(x,y)\mid X_u(x)=Y_t(y)=1]
       =G(x,y)\omega_k(G(x,y)).
\]
For $\mathcal B_a^{(j)}$, defined by replacing $k_a$ with $k_a^{(j)}$,
the right side is multiplied by $s_j(G(x,y)/f)$.
Apart from its start membership, its true value uses only the terminal
sampling matrix.
\end{lemma}

\begin{proof}
The row cutoff can fail to be full only if at least $K$ sampled columns
have $G(x,z)/a>r$, where $r=(1/2)/1000$ is fixed.
There are at most $M_G/(ar)$ such deterministic columns.
Conditional on $Y_t(y)=1$, the specified column contributes at most one
hit and each other column has $Y_a$-hit probability at most $C a$, by
Lemma~\ref{lem:env-ranks}.  Thus the expected count is at most
$1+CM_G/r$.  Conditioning also on the independent source sampling matrix does
not alter this bound.  Markov's inequality makes the failure probability
at most $1/8$ by a sufficiently large $K$.

Where $k_a(\mathsf{d}_y)>0$, we have $h(\mathsf{d}_y)=1$.  Therefore
$q_y=\mathbb E[R\mid Y_t(y)=1]\in[7/8,1]$ there, and the clipping
does nothing.  The conditional mean of Equation~\eqref{eq:atom-dense-table}
is consequently $k_a(\mathsf{d}_y)$.  Outside that support both sides vanish.
If a row output is positive, its first entry exceeds $1/2$ and its
$K$th largest first-coordinate value is below $1/1000$.  Hence there are at most $K-1$
positive entries.  Boundedness follows from the clipped denominator.  The routed
fraction is deterministic at truth and obeys the same argument.
\end{proof}

For the exact all-bin identity, the same true-table formulas also have a
meaning at scales with $P2^{-k}<1$. Here $L_k=1$, and every narrower usage box is the same
singleton $\{0\}^H$. Distinct identifiers have distinct sampling matrix
evaluations: writing $r=h_0-1$ and
$\bar v(y)=(y,\ldots,y^r)^{\mathsf T}$, we have
\[
 a(y)-a(z)=F\bigl(\bar v(y)-\bar v(z)\bigr)\ne0
 \qquad(y\ne z),
\]
because $F$ has independent columns and the first coordinate of
$\bar v(y)-\bar v(z)$ is $y-z\ne0$ in $\mathbb F_P$.
Thus at most one entry in a row is retained. Since $K>1$, its row
cutoff is identically one. On the support of $k_a(\mathsf{d}_y)$ we also have
$h(\mathsf{d}_y)=1$, so $q_y=c(q_y)=1$ and the bin table identity holds at
these finer scales as well, including its routed form. The comparisons
$\pi_a\asymp a$ are needed only on the admissible scales.
For each fixed positive entry, only finitely many $\omega_k$ are
nonzero; summing these true identities therefore recovers
Equation~\eqref{eq:atom-bin-partition} exactly. Section~\ref{sec:est}
will estimate the omitted deterministic tail without sampling it.

\subsection{Corrections with an incoming cutoff}

Corrections also need a bound on incoming degree.  Their deterministic
column-support bound permits an additional incoming cutoff while keeping
the probability that both cutoffs are full bounded away from zero.

Consider $F=fQ(G/f)g_0(y)$ and put $a=f$, $\vartheta=\theta_1$.
Fix $B_Q>\max\{K_G,\theta_2\}$. In the table formula, evaluate $Q$
at $\operatorname{clip}_{[0,B_Q]}(d)$. This bounded function has a
Lipschitz $s$th root on the real line, is zero for negative arguments,
and agrees with $Q(d)$ for every true argument $d=G(x,y)/f$.
The exact hierarchy continues to use the uncapped $Q$.  Let $R=R_{\theta_1}(\rho_x)$, $h=h_{\theta_1}$, and
$I=I_{\theta_1}$ be the functions in
Equation~\eqref{eq:atom-threshold-template}, where the true row order statistic is
$\rho_x=\operatorname{largest}_K\{G(x,z)/f:Y_f(z)=1\}$.
On a true or computed raw table, form $I(x,y)$ from that table's own row
entry and own order statistic.
With $\chi_{1/4,1/2}$ as above, define
\begin{equation}
 T_y=1-\chi_{1/4,1/2}\left(
       \operatorname{largest}_{K_c}\{I(x,y):X_u(x)=1\}\right).
 \label{eq:atom-incoming-gate}
\end{equation}
Thus $T_y=1$ through order statistic $1/4$ and $T_y=0$ from order statistic $1/2$.
This order statistic is taken over the fixed retained source set before any incoming
port filtering; there is no cutoff-dependent change of its index set.
The true correction preparation is $h(\mathsf{d}_y)R T_y$ on $Y_a$.
Define its conditional mean by
\begin{equation}
 q_y=\mathbb E\bigl[h(\mathsf{d}_y)R T_y\mid X_u(x)=Y_t(y)=1\bigr].
 \label{eq:atom-correction-denominator}
\end{equation}
The number $q_y$ is deterministic; as before, its second raw signal
carries the $Y_t$ membership mask.  Define the normalized true correction
using the clipped expression everywhere:
\begin{equation}
 \mathcal F(x,y)=\mathbf1_{\{Y_t(y)=1\}}
     \frac{Q(\operatorname{clip}_{[0,B_Q]}(\mathsf{d}_y))R T_y}{c(q_y)}\,g_0(y).
 \label{eq:atom-correction-table}
\end{equation}
Here $q_y=0$ whenever $\mathsf{d}_y\le\theta_1/2$.  Thus both kinds of second
signal are supported above a fixed positive first-signal threshold.
Choose $\zeta\le10^{-6}\min\{\theta_1,\delta_0,1/2\}$, decreasing it
later if a fixed compression margin requires this.  The second signals
are then uniformly tiny on the first-signal threshold scale.

\begin{lemma}[Correction table identity, including zero multipliers]
\label{lem:atom-correction}
For sufficiently large fixed $K$ and then $K_c$, the true correction
table satisfies
\[
 f\,\mathbb E[\mathcal F(x,y)\mid X_u(x)=Y_t(y)=1]=F(x,y).
\]
This holds also when $g_0(y)=0$, without a positive lower bound on
the denominator's unclipped mean. The cutoffs and normalization use only
the two role sampling matrices $X,Y$; $g_0$ is deterministic.
\end{lemma}

\begin{proof}
Write $Q_{\rm at}(d)=Q(\operatorname{clip}_{[0,B_Q]}(d))$ within this
proof. The row count argument in Lemma~\ref{lem:atom-dense}, with the correction
thresholds, gives $\Pr(R\ne1\mid\text{the two hits})\le1/8$.
If $g_0(y)>0$, every source with $I(x',y)>0$ belongs to the deterministic
set in Equation~\eqref{eq:atom-column-input}.  Conditional on $X_u(x)=1$,
the number of sampled members of that set has expectation at most
\[
 1+C uR_*/f\le1+C R_*.
\]
The bound remains valid after conditioning on the independent terminal
hit; it ignores the auxiliary row cutoff and therefore needs no additional
independence between rows.  If fewer than $K_c$ such sources are sampled,
the $K_c$th largest auxiliary score is zero and $T_y=1$.  Choose
$K_c\ge8(1+C R_*)$ to make $\Pr(T_y\ne1\mid\text{the two hits})\le1/8$.
The union bound gives
$\mathbb E[RT_y\mid\text{the two hits}]\ge3/4$.
Since $h(\mathsf{d}_y)=1$ on $Q$'s support, at every position with
$Q(\mathsf{d}_y)>0$ and $g_0(y)>0$ we have
\begin{equation}
 q_y=\alpha(x,y):=
 \mathbb E[RT_y\mid X_u(x)=Y_t(y)=1]\ge3/4.
 \label{eq:atom-alpha-bound}
\end{equation}
Thus $c(q_y)=q_y$ there. No lower bound on $q_y$ or $\alpha(x,y)$
is required when $g_0(y)=0$.  The claimed conditional
identity follows by division by this exact conditional mean.  If
$Q(\mathsf{d}_y)=0$ or $g_0(y)=0$, both sides of the identity are zero.  In the latter
case the pre-incoming quantity $Q_{\rm at}(\mathsf{d}_y)R/c(q_y)$ is still a well-defined
bounded comparison target, even if the unclipped mean vanishes.
\end{proof}

\subsection{Detector scores}

The unnormalized detector table is
\begin{equation}
 \mathcal S_i(x,y)=\phi(G_i(x,y)/f)
 R_{\delta_0}\!\left(\operatorname{largest}_K
       \{G_i(x,w)/f:Y_f(w)=1\}\right)
 \quad\text{on }X_f(x)=Y_f(y)=1.
 \label{eq:atom-detector-table}
\end{equation}
Its column trial is the $J$th largest such score over retained starts.
By the identical cutoff choice, its conditional mean given $Y_f(y)=1$
is exactly $U_i(y)$ from Equation~\eqref{eq:hier-detector}.
The trial is an inverse order statistic access to a row-bounded table and does not
require an incoming degree bound.

\subsection{Actual ports and estimated assembly}

The conditional-mean identities now give the true contributions.  A
product estimate will query a correction at $(x,z)$ and then a bin at
$(z,y)$.  Bounded outgoing degree limits the possible midpoints; bounded
incoming degree is needed when errors at the new source $z$ are summed
over old sources $x$.  We therefore specify access relations as well as
nonzero table supports.  The access relations may contain entries of
weight zero.

Passing through one copy projection is partitioned by its fixed source
and destination symbols $i,j$.  A task's \emph{inherited suffix} on a
sampling matrix is the string that the projections through which it was called
have appended there.  If the stage-$(l+1)$ copy sampling matrices have suffixes
$\xi_X,\xi_Y$, the stage-$l$ parent sampling matrices use $i\xi_X,j\xi_Y$.
Equation~\eqref{eq:env-suffix-identity} identifies their evaluations
pointwise with the copy evaluations, including the source, broad terminal
and narrower usage tests.  Separately, Lemma~\ref{lem:env-ranks}(4)
gives the unchanged joint law under each common suffix, so the stage-$l$
identities and sampling estimates apply to the parent table.  The task
tuple records these fixed symbols and suffixes.  Neither a right factor's
task parameters nor a column vector's task parameters retain an earlier
source address.
The environment supplied to the table is common to all its entries.
The vertices $x,y$ are arguments of that table, rather than additional
parameters.  Thus two paths that discover the same entry in that
environment query the same value; the path used to discover it is not
part of its mathematical definition.

We will apply these row functions to arrays represented by a finite list
of exceptional values and one common value at every omitted endpoint.
Call that common value the \emph{default}.  Adjoin arbitrarily many dummy
coordinates with the same default, and include them in every order statistic; they
are never matrix entries or ports.  For the true row arrays, dummy values
are zero.  Section~\ref{sec:buf-compression} constructs the finite-list
representations.  The elementary order statistic convention below lets us prove
support and access bounds before choosing that construction.

\begin{remark}[Order statistics of arrays with a common default]
\label{rem:atom-finite-ranks}
Suppose a row array equals a common default outside a finite list,
and arbitrarily many dummy coordinates have that same default.  Its
$k$th largest value can be computed from the listed coordinates and $k$ copies of
the default, with the stipulated zero padding.  Indeed, every omitted
coordinate merely supplies another copy of a value already available
with multiplicity at least $k$.  Additional listed real coordinates
equal to the default therefore do not change the order statistic.  The same
observation applies to an order statistic of entry sizes when each omitted entry
has the same default vector.

This is an equality of order statistics for one prescribed full array.  The dummy
coordinates are never matrix entries or ports.  A common componentwise
clamp of the entries and their default preserves this representation:
an omitted position still equals the clamped default.
\end{remark}

For the following access classes, $v=(v_y)$ is the full first raw
row, including its dummy coordinates, and $\rho=\operatorname{largest}_K(v)$.
For each fixed $\vartheta>0$, use repeatable approximations
$\widetilde v_y,\widetilde\rho$ with absolute error at most
$\eta=\vartheta/10000$. Define the classes of real retained endpoints by
\begin{equation}
 A=\{y:Y_a(y)=1,\quad
       \widetilde v_y>3\vartheta/8,\quad
       \widetilde\rho<\vartheta/200\}.
 \label{eq:atom-exposed-class}
\end{equation}
For access to the auxiliary score, define
\[
 A_I=\{y:Y_a(y)=1,\quad
       \widetilde v_y>\vartheta/10,\quad
       \widetilde\rho<\vartheta/40\},
\]
and use the smooth function
\[
 J_{\vartheta}(v,\rho)=
  \chi_{82\vartheta/1000,\,90\vartheta/1000}(v)
  \bigl(1-\chi_{26\vartheta/1000,\,27\vartheta/1000}(\rho)\bigr).
\]
The scalar functions in these formulas are evaluated entrywise.

\begin{lemma}[Support margins and exposed row classes]
\label{lem:atom-margins}
Suppose a bounded nonnegative main function $k(v)$ vanishes for
$v\le\vartheta$.  Its product $k(v)R_{\vartheta}(\rho)$ and the
preparatory product $h_{\vartheta}(v)R_{\vartheta}(\rho)$ have at most
$K-1$ nonzero real entries in a row.  Their common exposed class
$A$ has at most $K-1$ real entries, contains their supports, and satisfies
$I_{\vartheta}=1$ throughout $A$.  These assertions hold for arbitrary
real raw estimates, without a sign or accuracy hypothesis.

The auxiliary score $I_{\vartheta}$ also has at most $K-1$ nonzero real
entries. Its support is contained in $A_I$, which has at most $K-1$
entries, and $J_{\vartheta}$ is a bounded-support smooth witness equal
to one throughout $A_I$.
\end{lemma}

\begin{proof}
For the main and preparatory products a nonzero entry has
$v>\vartheta/2$ and $\rho<\vartheta/500$.  There cannot be $K$ such
entries, since then the $K$th largest value would exceed $\vartheta/2$.

Every positive main or preparatory entry belongs to $A$.  An entry in
$A$ satisfies $v_y>(3/8-1/10000)\vartheta>\vartheta/4$ and
$\rho<(1/200+1/10000)\vartheta<\vartheta/100$, so its witness
$I_{\vartheta}$ equals one.  Its entry threshold exceeds its order statistic
threshold, proving $|A|\le K-1$.

The support of $I_{\vartheta}$ requires $v>\vartheta/8$ and
$\rho<\vartheta/50$, again giving the same row bound.  The same strict threshold separation shows that $A_I$
contains this support and has at most $K-1$ entries.
The approximation margins imply $A_I\subseteq\{J_{\vartheta}=1\}$.
Moreover $82/1000>3(27/1000)$, so this witness also has bounded row
support and the required entry versus order statistic separation.  The witness is
an analytical output.

These arguments use the order statistic of the prescribed entire raw array.
In particular, if at least $K$ dummy positions have a common virtual
default $\beta$, then $\rho\ge\beta$.  A real position still equal to
that default cannot meet either exposed class: the required lower bound
on $v=\beta$ exceeds the required upper bound on $\rho$.  Thus omitted
default positions do not create unlisted ports.
\end{proof}

The fixed-accuracy approximations in this lemma are part of the access
interface, not tests for equality of exact reals.  The implementation
must use the same approximation and tie conventions whenever it queries
the same table.  The strict margins prove correctness for every such
choice.  Auxiliary-only access classes require a row bound but no
incoming bound.
In particular, different computed versions may expose different
zero-weight entries.  The assertions are applied to each version's own
final raw array and repeatable approximation convention.  Accuracy of the raw array is used
later in the error estimates, not in this support argument.


An estimated table applies the preceding bounded row functions to its
\emph{final} raw row arrays.  In a normalized task its second input is a
computed value approximating $\zeta\mathbf1_{\{Y_t(y)=1\}}q_y$: the
calculator divides that input by $\zeta$ and applies $c$.  It does not launch another averaging
operation after selecting an endpoint.  Its auxiliary incoming scores
are formed from its own final first-signal rows, and its incoming cutoff
is Equation~\eqref{eq:atom-incoming-gate} for those scores.  A preparatory
task is an earlier, separate task; its raw arrays, scores, and cutoffs are
its own, and its target is the unnormalized preparatory table specified
above.  Corrections additionally multiply by a clipped vector estimate
of the prescribed $U$-cutoff $g_0(y)$, queried with a common broad
$Y_f$-domain tuple.  The unnormalized preparations $hR$ and $hRT$
remain on their broad $Y_a$ domain; their usage parameter $t$ specifies
the conditional slice on which they are averaged.  The resulting
second raw signal and the normalized main output carry the $Y_t$ mask.
Bins and their preparations and detector tables have
no additional incoming multiplier.
For example, the root of a normalized pre-incoming function is,
on its fixed usage domain,
\[
 k(v)^{1/s}R(\rho)^{1/s}c(w/\zeta)^{-1/s}.
\]
The first two factors have the root regularity already proved, and
$q\mapsto q^{-1/s}$ has bounded derivative on $[1/2,1]$.
Since $c$ is $1$-Lipschitz, changing $w$ changes the last factor
by at most $C\zeta^{-1}|w-w'|$. Splitting a difference of the
three-factor product one factor at a time proves a global Lipschitz
bound on the saturated inputs. The same argument applies to the
unnormalized preparations and the score and witness products.
The constants may depend on the fixed $\zeta$, but not on the stage,
rates, or budget. Usage masks are the same in the compared tuple
and multiply both sides of these inequalities.

For clarity, Table~\ref{tab:atom-compression-thresholds} records the
thresholds needed by the compression argument.  For a pre-incoming row
function $H$, its entry factor vanishes through $e_H$, and its row cutoff
vanishes from order statistic $g_H$ onward.  Every listed pair satisfies
$3g_H<e_H$, including the broadest witness, for which
$82\vartheta/1000-3(27\vartheta/1000)=\vartheta/1000>0$.
All cutoffs are full in a neighborhood of order statistic zero.  The correction
preparation in this table is its pre-incoming factor $hR$; its incoming
cutoff is assembled afterward.

\begin{table}[htbp]
\centering
\begin{tabular}{@{}lll@{}}
\toprule
Pre-incoming row function & $e_H$ & $g_H$ \\
\midrule
Bin, including routed bins & $1/2$ & $1/1000$ \\
Bin preparation $hR$ & $1/4$ & $1/1000$ \\
Correction clipped-ratio function & $\theta_1$ & $\theta_1/500$ \\
Correction preparation $hR$ & $\theta_1/2$ & $\theta_1/500$ \\
Detector score & $\delta_0$ & $\delta_0/500$ \\
Auxiliary score $I_{\vartheta}$ & $\vartheta/8$ & $\vartheta/50$ \\
Witness $J_{\vartheta}$ & $82\vartheta/1000$ & $27\vartheta/1000$ \\
\bottomrule
\end{tabular}
\caption{Entry and row-cutoff thresholds. The entry factor vanishes through
$e_H$, and the row cutoff vanishes from order statistic $g_H$ onward. Here
$\vartheta\in\{1/2,\theta_1,\delta_0\}$ as appropriate.}
\label{tab:atom-compression-thresholds}
\end{table}

The bounded first-entry ranges and cutoff transition intervals already
specified give fixed sensitivity ranges for these functions.  If $w$
is a computed second raw entry, its only use is through $c(w/\zeta)$;
clipping $w$ to $[\zeta/2,\zeta]$ leaves this quantity unchanged.
Its denominator stays in $[1/2,1]$ even for negative or inaccurate
inputs.  Usage masks only remove entries and do not change any of these
support or continuity bounds.

For a correction or its preparation, let $A$ be the specified
pre-incoming class in Equation~\eqref{eq:atom-exposed-class}.
For a correction, the capped relation below is selected from this
uncapped row access class.
Define its capped relation by retaining $(x,y)\in A$ precisely when
\begin{equation}
 \#\{x':X_u(x')=1,\ (x',y)\in A\}\le K_c.
 \label{eq:atom-port-cap}
\end{equation}
This condition removes every entry at an overfull column, including
entries of weight zero.

\begin{lemma}[Port contract]
\label{lem:atom-ports}
All main and preparatory row tables, auxiliary scores, and the smooth
witnesses just constructed have bounded values and bounded row support.
Every row access class used for a main or preparatory output, or for
access to its auxiliary score, has bounded size and a smooth witness
equal to one on every exposed entry, including exposed entries of weight
zero.  A witness itself is needed as an analytical row function. The capped relation in Equation~\eqref{eq:atom-port-cap}
preserves every correction or preparatory table value and has at most
$K-1$ outgoing and $K_c$ incoming exposed entries.

These assertions hold separately for every computed version, provided
that its values, order statistics, scores, and access classes use the same final
raw arrays and common task tuple.  Access ports denote distinct entries:
the implementation must expose each pair at most once and pair its
incoming and outgoing ports consistently.  Access to auxiliary scores
alone has only the asserted row bound.
\end{lemma}

\begin{proof}
The row and witness assertions are Lemma~\ref{lem:atom-margins}; clipping
and multiplication by bounded factors do not enlarge supports.
If a column has more than $K_c$ entries in $A$, its auxiliary score is
one at all those entries, so its $K_c$th largest auxiliary score is one.
Equation~\eqref{eq:atom-incoming-gate} makes $T_y=0$ there.  Consequently
every value deleted by Equation~\eqref{eq:atom-port-cap} was already zero.
The surviving relation has the stated two degree bounds even when zero
weights are retained.  The argument applies to arbitrary raw arrays, so
it applies equally to the estimated and true versions.

The cap removes every uncapped row access port at an overfull column; it does
not choose $K_c$ of them.  It is applied even when the particular entry
has weight zero.  Every surviving port still lies in $A$, so its
pre-incoming witness $I$ equals one regardless of a zero usage mask,
incoming cutoff, or column multiplier.  The broader auxiliary class
$A_I$ is used to enumerate scores for an incoming order statistic.  It is not the
capped correction relation used as a left factor with a rotation map in a product.

The classes are relations on row and column IDs, rather than multisets
of paths.  Choosing one port per relation element gives the stated
combinatorial interface.  Establishing its single-address enumeration
and port pairing within the space bound is an implementation obligation;
the degree argument does not itself grant random access to both IDs.
\end{proof}

For later reference, the full interface supplied here is the following.
Each task has fixed common parameters; first-coordinate order statistics include the prescribed
dummy coordinates; exposed ports are real distinct entries and lie in a
unit witness; sparse main and preparatory ports obey the actual incoming
cap; and values, auxiliary scores, incoming order statistics, and column multipliers
are all queried at the same supplied environment and tuple.  The path
by which an entry is discovered must not change that environment or tuple
for its assembly.  Only completed earlier output types, or stated fixed local
functions and aggregations of those types, may be used as raw inputs of
a later task.  No prefix still waiting for assembly is such an output.
These requirements distinguish finite-support analytical bounds from
the access properties used in the subsequent implementation.
Completing a raw array here means defining its entries from the allowed
earlier outputs, rather than storing the whole array at once.  Its
auxiliary scores and uncapped row access classes are defined before its incoming
cutoff and cap, so their definitions do not call that same cutoff or cap.

These contracts apply directly to each computed array after compression
and clamping.  The later comparison in Lemma~\ref{lem:buf-clamping}
preserves relevant row-function values under its stated common-increment
hypothesis.  Each calculation uses its own consistent classes,
and Lemma~\ref{lem:atom-ports} gives their support and degree bounds
separately.

\subsection{Integrated losses and assembly bounds}

It remains to control how errors in the row functions, incoming scores,
and detector vectors affect an assembled table.  The incoming degree
bound lets us sum these errors by columns without multiplying by the
number of possible source vertices.

For a real row array $Z$ define
\begin{equation}
 \|Z\|_{u,a;s}^s=
 \frac1N\mathbb E_\sigma\sum_x\frac{X_u(x)}u
                 \sum_{y:Y_a(y)=1}|Z(x,y)|^s,
 \qquad
 \mathfrak d_{u,a}(H,H')=
   \|H^{1/s}-(H')^{1/s}\|_{u,a;s}.
 \label{eq:atom-row-norm}
\end{equation}
Here $\sigma$ has the product law of the sampling matrices and independent
auxiliary bits specified in Section~\ref{sec:env}.  All table values are
extended by zero outside their membership domain.  The denominator $u$
is the nominal rate, and actual-rate comparability in
Lemma~\ref{lem:env-ranks} absorbs the resulting absolute constants.
For a vertex array on a sampling matrix with rate $u$, write
\begin{equation}
 \|v\|_{u;s}^s=
    \frac1N\mathbb E_\sigma\sum_x\frac{X_u(x)}u|v(x)|^s;
 \label{eq:atom-vertex-norm}
\end{equation}
the indicated role sampling matrix replaces $X$ when necessary.  These are also
the comparison norms for two sample-dependent versions.  All required
pre-incoming functions and smooth support witnesses are included among
the row outputs whose discrepancies are to be bounded.

Bounded row support and range imply a uniform bound on these row-output
norms, since $|V_l|\le N$ and $\mathbb E X_u/u=O(1)$.  Bounded vertex
arrays have the analogous uniform bound.  Ordinary error is controlled
by root error on bounded nonnegative outputs, since
\begin{equation}
 |z-w|\le s C^{(s-1)/s}|z^{1/s}-w^{1/s}|
 \quad(0\le z,w\le C).
 \label{eq:atom-ordinary-from-root}
\end{equation}

\begin{lemma}[Assembly comparisons]
\label{lem:atom-assembly}
Compare two versions with the same environment, rates, and task tuple.
For a sparse normalized output write $H=P T g$ and
$H'=P'T'g'$, where $P,P'$ are the pre-incoming clipped-ratio row
functions (including their usage masks), $I,I'$ are their auxiliary
scores, $T,T'$ are the incoming cutoffs of
Equation~\eqref{eq:atom-incoming-gate} formed from $I,I'$, respectively,
and $g,g'$ are their column multipliers.  If $g,g'$ are the
prescribed smooth cutoffs of bounded vertex arrays $U_i,U_i'$, then
\begin{equation}
 \mathfrak d_{u,f}(H,H')^s
 \le C\mathfrak d_{u,f}(P,P')^s
    +C\mathfrak d_{u,f}(I,I')^s
    +C\frac fu\sum_i\|U_i-U_i'\|_{f;s}^s.
 \label{eq:atom-assembly-bound}
\end{equation}
For a sparse preparation omit the last term and use its unnormalized
pre-incoming row function.  The corresponding bin and detector row
functions have no incoming or column-multiplier assembly cost.

If $V_y,V_y'$ are the detector trials formed from
$\mathcal S_i,\mathcal S_i'$ by their column order statistic, then
\begin{equation}
 \|V-V'\|_{f;s}^s
   \le C\mathfrak d_{f,f}(\mathcal S_i,\mathcal S_i')^s.
 \label{eq:atom-detector-comparison}
\end{equation}
\end{lemma}

\begin{proof}
All factors are bounded.  On the union $E$ of the two final nonzero
supports, the triangle inequality for the products of $s$th roots gives
\[
 |H^{1/s}-(H')^{1/s}|
 \le C|P^{1/s}-(P')^{1/s}|
     +C|T^{1/s}-(T')^{1/s}|
     +C|g^{1/s}-(g')^{1/s}|.
\]
Outside $E$ the left side is zero.  The pre-incoming discrepancy may be
charged on its entire row domain.  Each column
has at most $2K_c$ edges in $E$, by Lemma~\ref{lem:atom-ports}.
By Equation~\eqref{eq:atom-rank-lipschitz}, Lipschitz continuity of the incoming
cutoff's $s$th root, and Equation~\eqref{eq:atom-ordinary-from-root},
\[
 |T_y^{1/s}-(T_y')^{1/s}|^s
 \le C\sup_{x:X_u(x)=1}
            |I(x,y)^{1/s}-I'(x,y)^{1/s}|^s.
\]
Consequently its integrated contribution is at most
\[
 \frac C{Nu}\mathbb E\sum_{y:Y_f(y)=1}
              \sum_{x:X_u(x)=1}
              |I(x,y)^{1/s}-I'(x,y)^{1/s}|^s,
\]
which is $C\mathfrak d_{u,f}(I,I')^s$.

The maximum over a fixed number of $U_i$ values is Lipschitz, and the
multiplier cutoff has Lipschitz root.  Thus the last discrepancy is
bounded in $s$th power by $C\sum_i|U_i(y)-U_i'(y)|^s$.
Summing over $E$ costs at most $2K_c/u$ per retained column.  The vertex
metric has weight $1/f$ instead, giving precisely the factor $f/u$ in
Equation~\eqref{eq:atom-assembly-bound}.  All columns in this comparison
are in the common $Y_f$ domain, including those whose final values carry
a narrower usage mask.  The common-tuple condition ensures that the
vertex discrepancy at $y$ is the same quantity for every old start.
Explicitly, if $m_E(y)$ counts the edges of $E$ entering $y$, then for
each $i$ its contribution obeys
\[
 \frac1{Nu}\mathbb E\sum_{y:Y_f(y)=1}
       m_E(y)|U_i(y)-U_i'(y)|^s
 \le 2K_c\frac fu\|U_i-U_i'\|_{f;s}^s.
\]
This is the column-counting step.

For a detector trial, Equation~\eqref{eq:atom-rank-lipschitz} bounds the
trial discrepancy by the supremum of the input score discrepancies over
retained starts.  Its $s$th power is at most the sum of their $s$th powers.
The factor $Y_f(y)/f$ in the vertex metric is exactly the factor obtained
by reversing the two sums in the row metric with $u=a=f$.
Equation~\eqref{eq:atom-ordinary-from-root} gives
Equation~\eqref{eq:atom-detector-comparison}.
\end{proof}

The required task types are therefore the bins $\mathcal B_b$ of
$C_l$, the routed bins $\mathcal B_b^{(j)}$ of $D_l$, detector scores
$\mathcal S_i$, and the pilot and actual correction tables $\mathcal F$,
together with one earlier unnormalized preparation for each normalized
type.  Their true conditional-mean identities are
Lemmas~\ref{lem:atom-dense} and~\ref{lem:atom-correction}.  Their computed
versions use the same bounded row functions and the assembly contract
above.  The next sections bound their discrepancies through the
additive schedule and implement their accesses without storing two
unrestricted vertex addresses.
